\documentclass[11pt]{article}
\usepackage[utf8]{inputenc}
\usepackage[T1]{fontenc}
\usepackage{lmodern}
\usepackage[margin=1in]{geometry}
\usepackage{amsmath,amssymb}
\usepackage{textcomp}
\usepackage[dvipsnames]{xcolor}
\usepackage{enumitem}
\usepackage{parskip}
\usepackage{titlesec}
\usepackage[colorlinks=true,linkcolor=blue,citecolor=blue,urlcolor=blue]{hyperref}

\titleformat{\section}{\large\bfseries}{\thesection}{1em}{}
\titleformat{\subsection}{\normalsize\bfseries}{\thesubsection}{1em}{}
\titleformat{\subsubsection}{\normalsize\bfseries\itshape}{}{0em}{}

\title{\textbf{Salience, Framing, or Prior?\\ Attributing the Asymmetric Decay of Soft and Hard In-Context Constraints}}
\author{}
\date{}

\newcommand{\ph}{\textit{[placeholder --- fill in after experiments are run]}}

\begin{document}
\maketitle
\vspace{-2em}

\begin{quote}
\small\textit{Assembled from the six confirmed stage outputs of an \texttt{auto-research-discussion} run (seed topic T3.2), post research-taste audit (\textcolor{ForestGreen}{\textbf{PASS}}), post per-topic review (\texttt{review\_2026-08.md}: one major and two moderate findings applied --- premise hedging, one added 2026 neighbor, RQ6 power escalation rule). Status: \textbf{specification / paper draft --- no experiments have been run.} All numbers below are pre-registered design targets, not results.}
\end{quote}

\begin{abstract}
Deployed LLM agents are steered by two tiers of in-context rule: soft organizational policies (tone, scope, business preferences) and hard safety norms (refusal of harmful requests), and a prior report --- whose primary source we have not been able to independently verify, and which we therefore treat as a hypothesis to re-estimate rather than a settled premise --- puts the former's behavioral decay at roughly eight times the latter's over long conversations. Accordingly, the design first treats the gap itself as a replication target: an unmanipulated arm must reproduce the asymmetry directionally before any mechanism claim is made. Whether such an asymmetry is fixable by changing where and how forcefully a policy is stated, or is instead baked into what alignment training already reinforced, has never been measured, because contextual salience, wording strength, and alignment-training reinforcement are naturally confounded in deployment --- hard norms are simultaneously stated more forcefully, placed more prominently, and trained in more deeply than soft policies. We isolate the three candidate mechanisms by intervention: a bank of content-matched soft/hard constraint pairs, audited on both authoring-time covariates and pilot-time task difficulty, is injected into distractor-padded long-horizon conversations while wording strength, context position, and repetition count are crossed orthogonally; alignment-prior strength is triangulated with an independent proxy battery rather than left unmeasured; and probe-induced measurement reactivity is calibrated against a between-conversation design before any mechanism claim is trusted. Fitting discrete-time survival models and decomposing the resulting decay-gap coefficient, we report what share of the asymmetry traces to presentation and what share survives even a maximally boosted soft-policy presentation --- deciding whether durable soft policies are a placement problem or a training problem.
\end{abstract}

\section{Introduction}

Long-running LLM agents are governed by two qualitatively different kinds of in-context rule. Hard safety norms --- refuse this request, never reveal that credential --- are reinforced deeply during alignment training and rarely need restating for a model to keep honoring them. Soft organizational policies --- stay within this product's scope, keep this tone, prefer this vendor --- exist only because a system prompt says so, and unlike a hard-safety failure, a soft-policy violation rarely trips any alarm: an agent that quietly answers a question outside its declared role looks, on the surface, exactly like one behaving correctly. That silence is precisely what makes the failure dangerous --- an agent already drifting out of its intended business behavior can operate for many turns before anyone notices.

A recent report puts a number on the asymmetry between these two tiers: soft-policy adherence decays roughly eight times faster than hard-safety adherence over the course of a long conversation. Two honesty notes bound how much weight this figure can carry. First, we have not been able to trace it to a verifiable primary source, so this paper treats it as a reported-but-unverified premise, not an established fact: the study's no-manipulation replication arm exists precisely to re-estimate whether --- and by how much --- the asymmetry holds in matched pairs, and every mechanism-attribution claim below is explicitly conditional on that arm replicating the gap directionally. Second, if the gap does not replicate, the contribution of this work becomes the non-replication finding itself, which would be equally decision-relevant: it would mean the widely assumed soft/hard asymmetry is an artifact of unmatched comparisons. Granting, for now, that some asymmetry exists --- nobody knows why. The gap could be a matter of salience --- soft policies are usually stated once and left to fade into a growing context window. It could be a matter of framing --- soft policies are typically phrased as preferences (``we'd like you to...'') rather than imperatives (``you must never...''). Or it could be that no amount of placement or wording can compensate for the fact that hard norms are baked into a model's weights by alignment training in a way a system-prompt sentence never will be. Which of these is true is not a matter of taste: if presentation is the whole story, the fix is a cheap, deployment-time change any policy author can make today; if the alignment prior dominates, no amount of prompt engineering will close the gap, and durable soft policies instead require training-time investment. Treating the two as interchangeable, which current practice implicitly does, risks applying exactly the wrong fix.

The nearest existing work approaches this gap from adjacent but distinct angles, and it is worth naming the two closest before surveying the rest: PrefEval already measures how long-context distraction erodes adherence to a user's soft preferences, and TimeToInconsistency already models multi-turn behavioral drift as a survival process --- but neither compares soft against hard constraint decay within one matched, causally interventional design, which is exactly the gap this paper closes. Beyond these two, OffTopicEval and SafeDialBench each instrument one side of the asymmetry in isolation --- the former measuring policy-scope adherence across query types, the latter measuring hard-safety consistency under adversarial pressure --- with no shared design to compare the two tiers' decay rates directly. SinkTrack presupposes that positional salience loss is the operative mechanism of instruction forgetting and builds a mitigation around that assumption, without testing whether constraint type or alignment prior modulates decay under the same positional disadvantage. And SafetyFewTokensDeep shows that alignment training reinforces safety refusals only a few tokens deep into a response --- a precise account of how the alignment prior is encoded for hard norms, but one that has never been extended to in-context-injected soft policies or to decay across conversation turns.

Turning this question into a measurable design surfaces its own difficulties. In real deployments, salience, framing, and alignment-training reinforcement move together --- hard norms are simultaneously stated more forcefully, placed more prominently, and trained in more deeply --- so any observational comparison of soft and hard constraints conflates all three; only intervention on content-matched pairs can pull them apart, and even then the alignment prior cannot be manipulated directly, only proxied. Constructing those matched pairs is itself nontrivial: soft policies are naturally graded (tone, scope, style) while hard norms are naturally binary (refuse or comply), so an unmatched pair would turn a claimed type effect into a disguised difficulty effect. The very act of measuring decay threatens to contaminate it: a probe that exercises a constraint's compliance necessarily touches its subject matter, and if that mention itself refreshes the constraint's salience, the measurement instrument is perturbing the quantity it is trying to observe --- differentially, perhaps, by constraint type. And because an unmanipulable mechanism (the alignment prior) is one of three candidate explanations, ruling nothing in or out for that mechanism would leave any attribution to it merely a diagnosis of exclusion rather than a measured finding.

Our approach turns each of these obstacles into a specific, checkable design choice rather than a hedge. A matched-pair bank is constructed and then audited twice --- once on authoring-time covariates (topic, length, verifier type) and once, after piloting, on realized task difficulty --- so a soft/hard contrast is trustworthy only after both checks pass. Wording strength, context position, and repetition count are then crossed orthogonally within one fitted hazard model, so a single specification yields both an average decomposition of how much of the gap presentation explains and a specific, pre-registered contrast at a maximally boosted soft-policy presentation, revealing how far presentation alone can go. The unmanipulable alignment prior is triangulated with three independent, imperfect proxies rather than left as an unexamined residual, so that whatever the design cannot manipulate, it can still measure. And before any of these results is trusted, the measurement instrument itself is calibrated: a between-conversation probing design that cannot suffer from salience refresh is compared against the main within-conversation design, and only a demonstrably type-uniform reactivity gap licenses treating the cheaper design's results as clean.

We evaluate this design through a staged pilot-then-scale-up plan, drawing constraint content from established rule-following and harmful-instruction corpora and distractor turns from a large real-conversation pool, with a no-injection floor and an unmanipulated replication arm as the necessary sanity checks before any manipulated-factor claim is asked to carry weight.

In summary, we contribute:
\begin{itemize}[leftmargin=1.4em]
\item An identification and formalization of the three-way confound underlying reported soft/hard in-context constraint decay asymmetries: salience, framing, and alignment-training prior move together in deployment, making any observational comparison uninterpretable without intervention.
\item A doubly-audited, content-matched soft/hard constraint-pair bank that distinguishes authoring-time covariate matching from pilot-time difficulty matching, so a type contrast is not secretly a difficulty contrast.
\item A decomposition of the soft/hard decay-gap coefficient within a single orthogonal manipulation design, yielding both an average attribution across presentation factors and a specific, pre-registered test of how far a maximally boosted presentation can close the gap.
\item A triangulated proxy battery for the unmanipulable alignment-prior mechanism, turning an otherwise-unexamined residual into a measured, checkable quantity.
\item A reactivity calibration of the measurement instrument itself, using a between-conversation design that cannot suffer from probe-induced salience refresh, run before any mechanism claim is trusted.
\end{itemize}

\section{Problem Definition}

\noindent\textbf{Formal setup.} A constraint $c$ has type $\tau(c) \in \{\text{soft}, \text{hard}\}$, wording variant $w(c) \in \{\text{advisory}, \text{standard}, \text{imperative}\}$, position $p \in \{\text{sys-head}, \text{sys-tail}, \text{mid-turn}\}$, and repetition count $r \in \{1, 3, 5\}$. A conversation runs turns $t = 1,\dots,T$ with the constraint statement injected per its $(w, p, r)$ assignment and probes scheduled at fixed intervals; probe outcome $A_c(t) \in \{0,1\}$ records whether the response at turn $t$ satisfies the constraint's verifier. A discrete-time hazard $h(t \mid \tau, w, p, r)$ is fitted by a mixed-effects model, $\text{hazard} \sim \text{turn} \times \tau \times w \times p \times r + (1\mid\text{pair}) + (1\mid\text{model})$, with survival $S(t)$ the corresponding cumulative non-violation probability. The soft/hard gap is $G = \exp(\beta_{\tau\times\text{turn}})$, the exponentiated type-by-turn coefficient; a decay half-life $t_{1/2}$ is the turn at which fitted survival first drops to half its turn-1 value. The object of study is the dominance-weight (partial-$R^2$) share of $\beta_{\tau\times\text{turn}}$ attributable to the manipulated factor blocks $\{w, p, r\}$ versus a residual share correlated with a proxy battery for $\tau$'s alignment-training prior, conditional on $G$'s confidence interval first excluding 1.

\noindent\textbf{Why this is hard.} Fixing a topic domain does not fix a matched soft/hard pair: soft constraints are naturally graded (tone, scope, style) while hard constraints are naturally binary (refuse or comply), so pairs matched only on domain and length can still differ in intrinsic difficulty, silently converting a $\tau$ effect into a difficulty effect. In ecological deployment, $w$, $p$, and the alignment-training prior are not independent of $\tau$ --- hard constraints are typically worded more imperatively, placed more prominently, and reinforced more heavily in training than soft ones --- so an observational estimate of $\beta_{\tau\times\text{turn}}$ confounds all three; only orthogonal manipulation of $(w, p, r)$ on content-matched pairs, holding $\tau$ fixed within a pair, can attribute variance to each factor separately. The alignment-training prior itself cannot enter as a manipulated factor at all --- it is fixed before the study begins --- so it can only be triangulated via proxies entered as covariates, yielding at best an upper bound on what correlational adjustment can claim about the residual. Finally, $A_c(t)$ is measured through probes that must exercise the constrained behavior to be checkable, and if a probe's mere occurrence refreshes $\tau$'s contextual salience, the hazard $h(t \mid \cdot)$ being fitted is itself an artifact of the measurement schedule rather than a property of the constraint --- a bias that must be measured and shown type-uniform before $\beta_{\tau\times\text{turn}}$'s decomposition can be trusted at all.

\section{Method Overview}

At a conceptual level, the study is a single measurement pipeline with four cooperating parts. A matched constraint-pair bank enters the pipeline pre-audited for both content-level comparability and realized difficulty; an orthogonal manipulation harness turns the resulting adherence sequence into a decomposed decay-gap coefficient and a specific boosted-presentation contrast; a proxy battery turns the same pairs' immovable alignment-training background into a measured covariate rather than an unexamined residual; and a reactivity calibration turns the whole measurement exercise into a claim about genuine constraint decay rather than an artifact of how often the instrument looks.

The \textbf{matched-pair bank} supplies content-matched soft/hard constraint statements, each authored in three wording-strength variants and passed through a two-stage audit --- covariate matching at authoring time, task-difficulty matching after a pilot run --- before entering any downstream analysis. The \textbf{orthogonal manipulation harness} crosses wording, position, and repetition within each constraint type and fits one hazard model whose type-by-turn coefficient is decomposed by dominance analysis across the manipulated factor blocks, while a specific boosted-presentation cell within the same fitted model is contrasted against standard-soft and standard-hard cells to test how far presentation alone can close the gap. The \textbf{prior-strength proxy battery} measures three independent covariates for the unmanipulable alignment-training mechanism and enters them into the same model, never as a manipulated factor. The \textbf{reactivity calibration} compares the main within-conversation probing design against a between-conversation design that structurally cannot suffer from probe-induced salience refresh, licensing (or withholding) trust in every result above.

\section{Methodology}

\subsection{Matched-Pair Bank Construction}
\textbf{Challenge.} Soft policies are naturally graded (tone, scope, style) while hard norms are naturally binary (refuse or comply), so a soft/hard pair matched only on topic domain and length can still differ in intrinsic difficulty, turning a claimed type effect into a difficulty effect in disguise.

\textbf{Insight.} Covariate matching alone cannot guarantee equal task difficulty, so treating realized turn-1 adherence parity as a second, empirically-checked matching criterion --- not an assumption --- is what keeps the eventual type contrast interpretable as a type effect.

\textbf{Method.} Soft/hard constraint pairs are authored matched on topic domain, token length (within 10\%), and verifier class (rule-based or judge-based), each in three wording-strength variants validated for semantic equivalence via embedding similarity ($\geq 0.90$) plus a human audit. A pilot run then checks each pair's turn-1 baseline adherence under standard wording and no manipulation; pairs whose soft and hard members differ by more than five points are dropped from the analysis bank rather than retained and hoped past.

\subsection{Orthogonal Manipulation Harness}
\textbf{Challenge.} Salience, framing, and alignment-training reinforcement move together in ecological deployment --- hard norms are stated more forcefully, placed more prominently, and trained in more deeply than soft ones --- so an observational comparison of soft and hard constraints cannot separate any one candidate mechanism from the other two.

\textbf{Insight.} A single fitted hazard model licenses two distinct, separately falsifiable readouts without being two mechanisms: an average partial-$R^2$ decomposition of the type-by-turn gap across manipulated factor blocks tests whether presentation matters on average, while a specific, pre-registered contrast at one maximally boosted presentation cell tests the practically decisive question of how far boosting actually gets.

\textbf{Method.} Wording strength, context position, and repetition count are crossed orthogonally within each constraint type, with total context token count held constant across cells via length-matched filler. A mixed-effects discrete-time hazard model is fit on the resulting adherence sequences, and the type-by-turn coefficient's variance is decomposed by dominance analysis across the three manipulated factor blocks. Within the same fitted model, a specific ``boosted-soft'' cell (imperative wording, system-head position, repetition five) is contrasted against the standard-soft and standard-hard cells to compute the fraction of the log-half-life gap it recovers.

\subsection{Prior-Strength Proxy Battery}
\textbf{Challenge.} The alignment-training prior cannot be intervened on at inference time --- it is fixed in a constraint's associated model weights before the study begins --- so without a credible proxy, this third candidate mechanism is un-investigable, and any residual left over from the manipulated factors is merely unexplained variance, not a measured finding.

\textbf{Insight.} Because the prior cannot be manipulated directly, triangulating it with multiple imperfect, differently-biased instruments and checking that they agree is what distinguishes a real residual-mechanism signal from proxy noise being mistaken for one.

\textbf{Method.} Three covariates are measured per constraint --- a no-injection zero-shot compliance rate, a base-versus-instruct compliance delta on an open-weights model pair, and, where model weights are accessible, a per-token KL divergence on the constraint-violating continuation --- and entered into the hazard model as covariates, never as manipulated factors.

\subsection{Reactivity Calibration}
\textbf{Challenge.} A probe that exercises a constraint's compliance must touch its subject matter to be checkable, and if that touch itself refreshes the constraint's contextual salience, the measurement instrument is perturbing the very decay it is trying to observe --- possibly to different degrees for soft versus hard constraints.

\textbf{Insight.} An anti-reactivity probe-authoring rule reduces but cannot provably eliminate salience refresh; the only way to bound the residual is to measure it directly against a design that structurally cannot suffer from it, and to check specifically whether that residual is uniform across constraint types.

\textbf{Method.} A calibration subsample of pilot-flagged cells is additionally run under a between-conversation design, in which each conversation is probed exactly once at a randomly assigned turn, so no earlier probe can refresh salience within that conversation. The fitted hazard curves from the between-conversation design are compared against the matched within-conversation design via an equivalence test on the probing-regime-by-type interaction.

\section{Experiments}

\subsection{Experimental Setup}

\textbf{Datasets.} LMSYS-Chat-1M supplies the distractor-turn pool used to pad conversations to a realistic long horizon without hand-authoring filler --- the same padding protocol PrefEval already validated for this purpose. IFEval supplies rule-checkable instruction-constraint templates for the subset of the pair bank that must be scored by a deterministic verifier rather than a graded judge, which the difficulty-artifact check below depends on. HEx-PHI and AdvBench supply established harmful-instruction taxonomies used to instantiate the hard-norm side of matched pairs with categories whose severity is already externally calibrated, rather than authored ad hoc.

\textbf{Baselines} (sanity-check comparators, not competitors to beat --- this is a measurement study, and no research question here is framed as ``our procedure beats baseline X''). A no-injection floor arm, in which the constraint statement is never injected, establishes that turn-1 adherence is caused by injection at all, so that ``decay'' is measured relative to a known floor rather than an undefined starting point. A no-manipulation natural-wording replication arm (standard wording, system-head position, repetition one) checks whether the reported soft/hard asymmetry replicates in matched pairs at all, independent of any manipulated-factor claim. A naive every-$k$-turn verbatim-restatement heuristic, already used informally in practice with no theoretical grounding, checks whether this design's own repetition-factor findings are at least consistent with existing ad hoc practice. SinkTrack-style training-free context anchoring, evaluated on open-weights models only, is included as the one prior method that treats positional salience as the operative mechanism of instruction forgetting, offering an informative comparison point for whether a competing salience-only intervention produces a similar or different effect size than this design's own manipulation. Finally, a strongest-alternative-explanation control (B-alt) targets the rival explanation that the gap is a difficulty/verifiability artifact --- that soft constraints are intrinsically harder to satisfy and scored by graded verifiers: analysis is restricted to pairs with rule-based binary verifiers on both sides, turn-1 adherence at or above 0.95 for both types, and no-injection floors within 10 points of each other, with the pre-registered check that the subset gap $\log G$ must lie within $\pm 50\%$ of the full-bank value; if the gap vanishes in this subset, the difficulty-artifact explanation wins and the mechanism-attribution claims (RQ2--RQ4) are withdrawn in favor of the artifact finding.

\textbf{Execution order.} RQ6 (reactivity calibration) must be evaluated first, on the calibration subsample, before RQ2 (factor significance), RQ3 (manipulated share), or RQ4 (boosted-presentation closure) are treated as trustworthy. If RQ6's equivalence test fails --- the probing-regime-by-type interaction's confidence interval extends outside the pre-specified equivalence bound --- the three gated research questions must be rerun on the reactivity-free between-conversation design, or their scope explicitly narrowed, before any mechanism-attribution claim is reported. This ordering is restated inside RQ6's own subsection below as a one-line reminder.

\textbf{Metrics.} The soft/hard hazard ratio $G$ and its confidence interval; decay half-lives per condition; the dominance-weight/partial-$R^2$ share of each manipulated factor block and of the prior-proxy covariates in the type-by-turn coefficient; the boosted-presentation gap-closure fraction $\kappa_{\text{close}}$; standard deviation and cross-correlation of the prior-strength proxies; and the probing-regime-by-type interaction hazard ratio from the reactivity calibration.

\textbf{Research questions:}
\begin{enumerate}[leftmargin=1.6em]
\item RQ1 --- Does the matched-pair construction procedure survive its own two-stage audit?
\item RQ2 --- Does at least one manipulated presentation factor significantly interact with the soft/hard decay gap?
\item RQ3 --- Do the manipulated presentation factors jointly explain the majority of the decay-gap coefficient?
\item RQ4 --- Does a maximally boosted soft-policy presentation close the gap with a hard norm's typical durability?
\item RQ5 --- Do the alignment-prior proxies carry shared, non-degenerate signal?
\item RQ6 --- Is the measurement instrument's own probe reactivity uniform across constraint types?
\end{enumerate}

\subsubsection{RQ1: Does the matched-pair construction procedure survive its own two-stage audit?}
\textbf{Setup.} As a first, minimum-viable check, 5 matched pairs restricted to rule-verifiable constraints (no judge-based verifiers) are authored and audited on covariate-matching criteria (topic, length, verifier class) at construction time, then piloted on one open-weights instruct model for 30-turn conversations with probes every 5 turns across 5 seeds (75 conversations total), checking each pair's turn-1 baseline-adherence difference. The full pilot bank (20 pairs, both rule- and judge-based) repeats the same two-stage audit before entering the main factorial.

\textbf{Expected outcome.} At least 80\% of constructed pairs survive both the covariate-matching construction criteria and the turn-1 baseline-adherence audit ($\leq$5-point difference) without requiring re-authoring.

\textbf{Falsification condition.} Refuted if fewer than 80\% of pairs survive both checks --- meaning either the matching procedure or the underlying premise that soft and hard constraints can be matched at all needs redesign before any downstream analysis is trustworthy.

\textbf{Hypothesis (plain language).} Soft and hard constraints can be authored into pairs that are comparable both on paper and in realized difficulty, at a survival rate high enough to support the main design without constant re-authoring.

\textbf{Metrics.} Pair-survival rate through each audit stage; distribution of turn-1 adherence differences among surviving pairs.

\textbf{Result.} \ph

\subsubsection{RQ2: Does at least one manipulated presentation factor significantly interact with the soft/hard decay gap?}
\textbf{Setup.} Pilot stage: 2 models $\times$ 20 matched pairs $\times$ a $2\times2\times2\times2$ factorial (type: soft/hard $\times$ wording: advisory/imperative $\times$ position: system-head/mid-turn-10 $\times$ repetition: 1/3), 30-turn conversations with probes every 5 turns, 3 seeds per cell, $\approx$960 conversations per model. Scale-up: $\geq$6 models across $\geq$3 model families $\times$ 40 pairs $\times$ 3-level wording and position factors $\times$ repetition $\{1, 3, 5\}$, horizon extended to 100 turns on a fractional-factorial subsample. Likelihood-ratio tests on the mixed-effects discrete-hazard model for each factor-by-type-by-turn interaction, Holm-Bonferroni corrected across the three mechanism-factor tests, $\alpha = 0.05$.

\textbf{Sample-size rationale.} The pilot's 20 pairs and $\approx$960 conversations per model are sized for $\geq$80\% power to detect an interaction hazard ratio of 1.5 at $\alpha = 0.05$ (to be confirmed by simulation before launch); an underpowered pilot result triggers the pre-committed scale-up rather than a null claim.

\textbf{Expected outcome.} At least one of wording, position, or repetition shows a statistically significant factor-by-type-by-turn interaction after correction.

\textbf{Falsification condition.} Refuted if none of the three manipulated factor blocks shows a significant interaction at corrected $\alpha = 0.05$ --- meaning presentation factors do not move the soft/hard gap at all within the tested design, and RQ3/RQ4's premise that presentation is worth decomposing further has no basis.

\textbf{Hypothesis (plain language).} At least one of how a soft policy is worded, where it is placed, or how often it is repeated measurably changes how fast it decays relative to a hard norm.

\textbf{Metrics.} Per-factor interaction coefficients and Holm-Bonferroni-corrected $p$-values; per-cell hazard ratios.

\textbf{Result.} \ph

\subsubsection{RQ3: Do the manipulated presentation factors jointly explain the majority of the decay-gap coefficient?}
\textbf{Setup.} Dominance analysis / partial-$R^2$ decomposition of the type-by-turn coefficient from the same fitted model as RQ2, comparing the combined share attributable to the wording, position, and repetition factor blocks against the share attributable to the prior-strength proxy covariates (RQ5).

\textbf{Expected outcome.} The manipulated factor blocks jointly absorb more than 50\% of the type-by-turn coefficient's partial-$R^2$.

\textbf{Falsification condition.} Refuted, toward ``presentation is not the story,'' if the manipulated share is 50\% or below while the prior-proxy covariates absorb the majority --- in which case matching a soft policy's presentation to a hard norm's still leaves most of the gap standing, and the residual traces to something presentation cannot reach.

\textbf{Hypothesis (plain language).} Most, but not necessarily all, of the soft/hard decay gap is explained by how constraints are presented rather than by an unreachable background difference.

\textbf{Metrics.} Dominance-weight/partial-$R^2$ share per factor block and for the prior-proxy covariates.

\textbf{Result.} \ph

\subsubsection{RQ4: Does a maximally boosted soft-policy presentation close the gap with a hard norm's typical durability?}
\textbf{Setup.} Within the same fitted model, the boosted-soft cell (imperative wording, system-head position, repetition five) is contrasted against the standard-soft and standard-hard cells; gap-closure fraction $\kappa_{\text{close}}$ and half-life ratios are computed with 10,000-resample bootstrap confidence intervals at the conversation level.

\textbf{Expected outcome.} $\kappa_{\text{close}} \geq 0.5$ with a bootstrap 95\% confidence-interval lower bound above 0.25, while the boosted-soft cell's half-life remains statistically distinguishable from the standard-hard cell's (half-life-ratio confidence interval excludes 1) --- presentation narrows the gap substantially without fully eliminating it.

\textbf{Falsification condition.} Refuted in either direction: toward ``presentation is inert'' if $\kappa_{\text{close}}$'s confidence interval includes 0; toward ``presentation is the whole story'' if the boosted-soft cell's half-life is statistically indistinguishable from, or exceeds, the standard-hard cell's.

\textbf{Hypothesis (plain language).} Maximally boosting how a soft policy is presented moves it substantially closer to hard-norm-level durability, but does not make it indistinguishable from a hard norm.

\textbf{Metrics.} $\kappa_{\text{close}}$ with bootstrap confidence interval; boosted-soft-versus-standard-hard half-life ratio with confidence interval.

\textbf{Result.} \ph

\subsubsection{RQ5: Do the alignment-prior proxies carry shared, non-degenerate signal?}
\textbf{Setup.} No-injection zero-shot compliance rate and base-versus-instruct compliance delta are measured across all constructed constraints on the open-weights model pair; their standard deviation and Spearman correlation with each other are computed.

\textbf{Expected outcome.} Standard deviation of the no-injection compliance rate exceeds 0.1, and its Spearman correlation with the base-versus-instruct delta exceeds 0.3.

\textbf{Falsification condition.} Refuted if the proxies are near-flat (standard deviation at or below 0.1) or uncorrelated/contradictory (Spearman correlation at or below 0.3) --- meaning the alignment-prior mechanism is uninstrumented, and any residual share attributed to ``prior'' in RQ3 must be reported as unidentified rather than prior-linked.

\textbf{Hypothesis (plain language).} The proxies for the unmanipulable alignment-training mechanism vary meaningfully across constraints and agree with each other, rather than being flat or contradictory.

\textbf{Metrics.} Standard deviation of each proxy; pairwise Spearman correlation between proxies.

\textbf{Result.} \ph

\subsubsection{RQ6: Is the measurement instrument's own probe reactivity uniform across constraint types? (evaluated first --- gates RQ2, RQ3, and RQ4)}
\textbf{Setup.} A calibration subsample (4 pilot-flagged cells $\times$ 2 constraint pairs $\times$ 6 probe turns $\times$ 10 seeds, $\approx$480 between-conversation runs) is probed exactly once per conversation at a randomly assigned turn and contrasted against matched within-conversation curves via a two-one-sided-tests (TOST) equivalence procedure on the probing-regime-by-type interaction hazard ratio, with equivalence bounds $[0.80, 1.25]$ and $\alpha = 0.05$.

\textbf{Power arithmetic and escalation rule.} A back-of-envelope calculation shows the 480-run floor cannot by itself decide the TOST: with per-cell event rates in the 30--50\% range, $\approx$120 conversations per regime-by-type margin yield a standard error for the interaction log-hazard-ratio of roughly 0.3--0.4, i.e.\ a 95\% confidence interval of multiplicative width $\approx \exp(\pm 0.6\text{--}0.8)$ --- wider than the $[0.80, 1.25]$ band (half-width $\log 1.25 \approx 0.22$) regardless of the true effect. The 480 runs are therefore pre-registered as a floor, not the decision sample, with a group-sequential escalation: the subsample is expanded in stages (pooling over probe turns, up to $\approx$10$\times$ the floor) until the interval's width alone no longer forces the outcome. The gate's failure semantics distinguish two cases accordingly: \emph{equivalence rejected} (interval excludes 1 and lies materially outside the band --- probe reactivity genuinely differs by type) triggers the between-conversation rerun of RQ2--RQ4; \emph{underpowered width} (interval centered near 1 but wide) triggers only the next escalation stage, never a spurious rerun of the expensive arm.

\textbf{Expected outcome.} The 95\% confidence interval of the probing-regime-by-type interaction hazard ratio lies entirely within $[0.80, 1.25]$.

\textbf{Falsification condition.} Refuted if the interaction's confidence interval extends outside $[0.80, 1.25]$ --- meaning probe reactivity differs by constraint type, biasing RQ2--RQ4's within-conversation estimates differentially rather than uniformly, and the main analyses must be rerun on the between-conversation design (or their scope narrowed) before any mechanism-attribution claim is reported. This is why RQ6 is evaluated before, not after, RQ2--RQ4.

\textbf{Hypothesis (plain language).} Probing the model to check compliance may refresh the constraint's salience somewhat, but not differentially enough between soft and hard constraints to bias the comparison the main design relies on.

\textbf{Metrics.} Probing-regime-by-type interaction hazard ratio with 95\% confidence interval; TOST equivalence decision.

\textbf{Result.} \ph

\vspace{0.5em}
\noindent\textbf{Further evaluation.} With more time, a natural extension is to test whether the attributed mechanism shares themselves vary across languages, tool-using agent settings, or multi-agent deployments --- explicitly out of scope here because each would introduce its own confound (translation-induced wording change; tool-call turns as a distinct kind of ``probe'') requiring its own matched design rather than a direct extension of the single-agent, English-only harness above.

\section{Related Work}

The closest prior work to this paper is PrefEval and TimeToInconsistency, and it is worth stating the delta plainly before surveying the wider landscape: PrefEval already measures how long-context distraction erodes adherence to a user's soft preferences, and TimeToInconsistency already models multi-turn behavioral drift as a survival process --- but neither compares soft against hard constraint decay within one matched, causally interventional design, which is precisely the gap this paper is built to close.

\textbf{Soft-policy and preference-adherence measurement.} PrefEval tests whether models infer, remember, and proactively follow a user's stated preferences across long, distractor-interleaved multi-session conversations; OffTopicEval instantiates purpose-specific agents via system-prompt policies and measures the accept/refuse tradeoff as query type shifts from in-domain to adversarially ``laundered'' out-of-domain requests. Both establish that soft, policy-like constraints are measurably fragile, but neither includes a hard-safety-norm arm to compare against, nor manipulates wording, position, or repetition factorially --- they document soft-constraint failure without attributing its cause.

\textbf{Multi-turn safety and dialogue-consistency evaluation.} SafeDialBench scores hard-safety consistency across multi-turn dialogues instrumented with seven jailbreak strategies, and ``LLMs Get Lost in Multi-Turn Conversation'' documents a large, general multi-turn performance drop decomposed into aptitude loss and a much larger rise in unreliability. Both establish that multi-turn interaction degrades LLM behavior broadly, but neither stratifies degradation by soft-versus-hard constraint type, and SafeDialBench's adversarial-only design leaves benign, attack-free decay of hard norms unmeasured --- the asymmetry this paper investigates is invisible to either design.

\textbf{Mechanistic accounts of context degradation.} SinkTrack diagnoses instruction forgetting as a consequence of recency bias and attention-sink anchoring, and proposes a training-free fix that restores access to early context; TimeToInconsistency models the first multi-turn inconsistency as a survival event with observational semantic-drift covariates; ICFBench isolates the causal effect of \emph{revoked} context via a paired design that suppresses outdated instructions. Closest to our salience factor block, the concurrent controlled study ``Prompt Design at Scale'' (arXiv:2607.19257) crosses instruction count, format, and system-prompt-versus-user-turn placement against context length and finds placement effects at least as large as format effects, with perfect adherence collapsing by 80 simultaneous rules --- direct interventional evidence that the presentation factors we manipulate are causally live; we concede that overlap plainly, and note that it manipulates neither wording strength nor constraint type, measures point-in-time adherence rather than fitted decay trajectories, and attributes no share of any adherence gap to alignment priors --- the type-asymmetry estimand this paper is built around is outside its design. Each names or models one candidate mechanism of context-related failure, but SinkTrack presupposes positional salience as the sole cause without testing whether constraint type or alignment prior also matters, TimeToInconsistency's covariates are observational rather than manipulated and so cannot separate salience from framing from prior, and ICFBench's paired design targets the opposite failure direction --- suppressing context that should be forgotten, not preserving context that should persist.

\textbf{Alignment-prior depth and wording/robustness diagnostics.} SafetyFewTokensDeep shows that safety alignment reshapes mostly the first few output tokens of a refusal, instrumented via per-token KL divergence between aligned and base models; TARE shows that prompt performance depends heavily on exact wording even among semantically similar paraphrases, formalizing this as worst-case loss over a paraphrase neighborhood; RECAST synthesizes multi-constraint instruction data using a dual-mode verifier that combines rule-based and LLM-judge checks. Each contributes an instrument this paper repurposes --- the aligned-versus-base divergence idea for the prior-proxy battery, the paraphrase-neighborhood discipline for validating wording-strength variants, the dual-mode verifier taxonomy for scoring compliance --- but none of the three measures decay of an in-context-injected constraint across conversation turns, which is the phenomenon this paper's design is built around.

\section{Conclusion}

This work reframes a reported but independently unverified number --- the roughly eightfold gap between soft-policy and hard-norm decay --- as a question with a precise identifiability structure, rather than a fact to be repeated: the design first re-estimates whether the asymmetry holds at all in matched pairs, and only then asks why. The central methodological wager is that salience, framing, and alignment prior are only separable once content-matched pairs are intervened on directly and the measurement instrument's own reactivity is calibrated before any attribution is trusted; a merely observational comparison, however large, cannot tell these three apart. Whichever share dominates, the finding is decision-relevant: if presentation accounts for most of the gap, policy authors have a cheap, immediate lever; if a substantial residual survives even maximally boosted presentation, that residual is evidence that durable soft-policy adherence requires training-time investment no prompt can substitute for. Two limitations bound this claim from the outset and are natural next steps rather than afterthoughts: the alignment-prior mechanism remains the weakest-instrumented of the three candidates, triangulated through proxies rather than directly manipulated, so any residual attributed to it is an upper bound on what correlational adjustment can claim, not a causal estimate; and the study is scoped to single-agent, English-only, sandboxable conversations, leaving multi-agent, tool-using, and cross-lingual deployments --- where salience, framing, and training-prior effects may combine differently --- for follow-up work building on the released matched-pair bank and calibration harness.

\vspace{1em}
\hrule
\vspace{1em}

\section*{Pre-handoff checklist}
\begin{itemize}[leftmargin=1.4em]
\item \textbf{Organization/navigation} --- narrative flows problem $\to$ method $\to$ experiments $\to$ related work without forward references to undefined terms.
\item \textbf{Terminology discipline} --- ``matched pair,'' ``decay-gap coefficient,'' ``boosted-soft cell,'' ``reactivity calibration'' each defined at first use in Problem Definition/Methodology, used consistently thereafter; no internal stage IDs (\texttt{sp\_*}, \texttt{m\_*}, \texttt{rq\_*}) leaked into prose --- checked directly against the discussion pipeline log.
\item \textbf{Notation with worked example} --- the formal setup section defines $(\tau, w, p, r, t, A_c(t), h, S, G, t_{1/2})$ and the Problem Definition/Methodology sections consistently reuse this notation; no separate worked numeric example is included since no numbers exist yet (would require fabricating results).
\item \textbf{RQ framing matches outcomes} --- all 6 RQs are graded/quantitative (rates, coefficients, ratios, proportions), each answered in its own subsection with an explicit falsifiable expected outcome and a distinct falsification condition, not a binary yes/no.
\item \textbf{No typesetting artifacts} --- only explicit placeholder markers in Results, per the drafting convention this is a pre-registered spec, not a report.
\item \textbf{Artifact plan} --- the matched-pair bank, the manipulation harness, and the reactivity-calibration subsample are the natural release artifacts; not yet built, so not claimed as available.
\end{itemize}

\section*{Hand-off}
This is a draft, not a finished paper --- experiments have not run. Natural next steps: prose polish and venue-specific LaTeX formatting once a target venue is chosen, or revising an earlier stage if piloting surfaces a problem (e.g., if the pair-bank survival rate in RQ1 is much lower than expected).

\end{document}
